\documentclass[11pt]{article}

\usepackage[margin=1in]{geometry}
\usepackage[T1]{fontenc}
\usepackage[utf8]{inputenc}
\usepackage{lmodern}
\usepackage{amsmath,amssymb,amsthm,bm,mathtools}
\usepackage{booktabs}
\usepackage{graphicx}
\usepackage{caption}
\usepackage{xcolor}
\usepackage{hyperref}
\usepackage{enumitem}
\usepackage{float}

\hypersetup{colorlinks=true, linkcolor=blue!60!black, citecolor=blue!60!black,
            urlcolor=blue!60!black}

\newtheorem{proposition}{Proposition}
\theoremstyle{remark}
\newtheorem{remark}{Remark}

\newcommand{\R}{\mathbb{R}}
\newcommand{\E}{\mathbb{E}}
\newcommand{\Var}{\operatorname{Var}}
\newcommand{\Cov}{\operatorname{Cov}}
\newcommand{\Xis}{\Xi}
\newcommand{\code}[1]{\texttt{\small #1}}

\title{\textbf{Layer 6: hierarchical priors, speciation, and memorization}\\[2mm]
\large A session report on belief propagation for diffusion scores}
\author{Session record --- branch \code{claude/em-bp-denoiser-learning-e07ike}}
\date{2 August 2026}

\begin{document}
\maketitle

\begin{abstract}
\noindent
This report documents one working session on the \code{GioviManto/diffusion}
project. The session was directed by a single instruction: study two papers and
take inspiration from them. Both turned out to be by the project's advisors.
The work that followed built a hierarchical (tree) prior with exact belief
propagation and exact EM on it, derived the two diffusion time scales in closed
form for that prior, and ran five experiments against them. It also produced two
corrections to my own earlier claims, one of them substantive, and one negative
result. Everything is on the feature branch; \code{main} is untouched. The test
suite went from 50 to 101 tests, all passing.
\end{abstract}

\tableofcontents

\section{Scope, and how to read this}

\subsection{What the session was asked to do}

The instruction was to study
\begin{itemize}[nosep]
  \item \textbf{P1} --- Garnier-Brun, M\'ezard, Moscato, Saglietti, \emph{How
    transformers learn structured data: insights from hierarchical filtering},
    arXiv:2408.15138; ICML 2025 (PMLR 267:18831--18847);
  \item \textbf{P2} --- Biroli, Bonnaire, de Bortoli, M\'ezard, \emph{Dynamical
    regimes of diffusion models}, Nature Communications \textbf{15}, 9957
    (2024); arXiv:2402.18491,
\end{itemize}
and continue the research. Marc M\'ezard is an author of both; J\'er\^ome
Garnier-Brun of the first.

\subsection{Provenance of the source material --- read this before citing}

This matters enough to state before any result.

\paragraph{P1 has been read in full.} Direct HTTP is blocked at this
environment's gateway: \code{arxiv.org}, \code{www.nature.com}, PMC, HAL and
PMLR all return \code{403} on \code{CONNECT}, from both the fetch tool and
\code{curl}. However the Hugging Face Hub exposes
\code{hf://papers/2408.15138/paper.md} through a server-side route, which
returns the complete \LaTeX ML rendering of the paper including all appendices.
Every statement about P1 in this report is from that source.

\paragraph{P2 has \emph{not} been read.} It is not available on that route and
every direct route is blocked. What is said about P2 here comes from web search
--- abstract-level statements about the two cross-overs and the entropic
criterion --- and \textbf{nothing mathematical is quoted from it}. The
speciation criterion used throughout is derived from scratch in
Section~\ref{sec:speciation} in this package's own conventions and checked
against sampled trajectories. One search summary rendered the condition with a
reciprocal, $t_S = \tfrac12\log(1/\Lambda)$, contradicting its own statement
that $t_S$ diverges with dimension; deriving it independently avoided importing
that slip. \textbf{P2 should be read before it is cited.}

\paragraph{A third paper matters and was not in the original list.}
Sclocchi, Favero and Wyart, \emph{A phase transition in diffusion models reveals
the hierarchical nature of data}, PNAS \textbf{122}(1) e2408799121 (2025),
arXiv:2402.16991 --- \emph{cited by P1} --- studies diffusion on hierarchically
generated data directly. It is the closest prior work to Section~\ref{sec:cascade}
and is discussed in Section~\ref{sec:novelty}. It has not been read in full
either; the account here is from search.

\subsection{Structure}

Section~\ref{sec:p1} summarises P1 as read. Section~\ref{sec:machinery}
describes the new code. Section~\ref{sec:theory} gives the derivations, all of
which are verified numerically. Section~\ref{sec:results} reports the
experiments. Section~\ref{sec:corrections} records what I got wrong and what
came out negative. Section~\ref{sec:limits} states the limitations and what
should happen next.

\section{P1, as read}
\label{sec:p1}

\subsection{The data model}

A root symbol $x_0$ is drawn from a vocabulary of size $q$. A transition
\emph{tensor} $\bm M \in \R_+^{q\times q\times q}$, with
$M_{abc} = P(\text{children } (b,c) \mid \text{parent } a)$ and
$\sum_{bc} M_{abc} = 1$, is applied for $\ell$ generations down a binary tree,
producing $2^\ell$ leaves. Entries are log-normally sampled and
\emph{non-overlapping}: if $M_{abc}>0$ then $M_{a'bc}=0$ for all $a'\neq a$. The
production rules are therefore unambiguous, and at full hierarchy the whole
tree including the root is deterministically recoverable from the leaves. The
authors describe it as a simplified probabilistic context-free grammar with a
fixed parse tree, and relate it to the Random Hierarchy Model of Cagnetta et
al.\ (uniform rates, layer-dependent rules); the differences are log-normal
versus uniform rates and a single tensor reused at every level.

\subsection{Hierarchical filtering}
\label{sec:filtering-def}

This is P1's central device and the part I originally implemented incorrectly
(Section~\ref{sec:correction-filtering}). At filter level $k$, the $2^k$ nodes
at depth $k$ are drawn \textbf{conditionally independently given the root},
each with the marginal the unfiltered model assigns it:
\begin{equation}
  P(x_j = b \mid x_0 = a)
  \;=\;
  \bigl(\bm p_0\, \bm M^{\sigma_0(j)} \bm M^{\sigma_1(j)} \cdots
        \bm M^{\sigma_{k-1}(j)}\bigr)_{ab},
\end{equation}
tracing the tensor along the unique root-to-$j$ path, with
$\sigma_m(j)\in\{L,R\}$ recording the branchings and $\bm M^{L},\bm M^{R}$ the
two marginalisations of $\bm M$. Levels $1,\dots,k-1$ are skipped. Below level
$k$ the ordinary process resumes.

The consequences are worth spelling out because they are exactly what makes
filtering a clean experimental knob:
\begin{itemize}[nosep]
  \item correlations survive inside blocks of $2^{\ell-k}$ leaves;
  \item blocks remain correlated \textbf{through the root} --- they are not
    independent;
  \item all marginals are preserved at every $k$;
  \item at $k=\ell$ the leaves are conditionally i.i.d.\ given the root, and P1
    notes that a Naive Bayes classifier is then optimal and every attention
    layer superfluous.
\end{itemize}

\subsection{Inference and what it is used for}

The graph is a tree at every $k$, so BP is exact and linear-time, converging in
$2(\ell-k+1)$ steps. This gives a \emph{family} of exact algorithms
$\mathrm{BP}_k$, and running $\mathrm{BP}_k$ on data generated at a different
filtering level is a well-defined \emph{mismatched oracle}. That family is P1's
measuring instrument, and adopting it is the main methodological import of this
session (Section~\ref{sec:exp15}).

\subsection{Findings}

\begin{enumerate}[nosep]
  \item Transformers match $\mathrm{BP}_{k_{\mathrm{train}}}$'s accuracy
    \emph{and its calibrated probabilities}, in-sample and out-of-sample,
    including where that oracle is not optimal --- evidence they implement the
    algorithm rather than merely scoring well.
  \item \textbf{Sequential acquisition.} During training the predictions align
    first with $\mathrm{BP}_\ell$ (leaf-to-root correlations only) and then with
    \emph{decreasing} $k$: a clean staircase, shortest-range correlations first.
  \item Attention maps develop blocks of size $\sim 2^{\ell-k}$, organised
    hierarchically across layers; layer-wise probes recover ancestors up to the
    depth of the probed layer.
  \item Appendix F gives an explicit construction embedding BP in $\ell$
    transformer layers, using $\mathcal O(q^2)$ memory slots per token to fold
    the downward pass into the upward one.
  \item MLM pre-training sharply reduces the labelled data needed for root
    classification.
\end{enumerate}

\subsection{P2, as far as search reports it}

The backward dynamics passes through two cross-overs. \emph{Speciation}, where
the coarse structure of the sample is decided by a symmetry-breaking mechanism,
with a time scale set by the top eigenvalue of the data covariance --- diverging
logarithmically when that eigenvalue grows with dimension. \emph{Collapse},
where trajectories are captured by a single training point through a mechanism
analogous to glassy condensation, with a time scale set by an excess entropy,
and the consequence that memorization is avoidable at finite times only if the
dataset is exponentially large in the dimension. The collapse statement is about
the \emph{empirical score}: the exact score of the measure placing mass $1/N$ on
each training point.

\section{What was built}
\label{sec:machinery}

Two new modules, three new experiments, 51 new tests. Everything is pure
\code{numpy}; no autodiff anywhere, consistent with the rest of the package.

\subsection{\code{src/hierarchy.py}}

\paragraph{The prior.} \code{GaussianTree(depth, branching, rho,
root\_separation, filter\_level)} is the continuous analogue of P1's model:
\begin{equation}
  z_{\text{root}} \sim \mathcal N(0,1), \qquad
  z_{\text{child}} = \rho\, z_{\text{parent}} + \sqrt{1-\rho^2}\;\varepsilon,
\end{equation}
with the leaves observed and internal nodes latent. Every node has unit marginal
variance, matching the normalisation the chain priors use, so tree and chain
families are directly comparable at a fixed noise schedule. For Gaussian
innovations the leaf covariance is ultrametric,
\begin{equation}
  \Cov(\text{leaf}_i,\text{leaf}_j) \;=\; \rho^{\,2(L - d_{ij})},
\end{equation}
where $d_{ij}$ is the depth of the lowest common ancestor.

\paragraph{Exact BP, twice.}
\begin{itemize}[nosep]
  \item \code{tree\_bp\_gaussian} --- information form $(h,\lambda)$, exact,
    $\mathcal O(n)$. Upward and downward updates touch only those two numbers
    per message, with nowhere for a projection step to hide. This is deliberate:
    the audit that preceded this work found the previous Gaussian BP was doing
    ``something strange'' instead of updating mean and variance.
  \item \code{tree\_bp\_grid} --- grid messages, valid for \emph{any} innovation
    law, using the same two primitives as the chain ($\bm K^\top f$ upward,
    $\bm K g$ downward) and the same \code{K[out, in]} convention. The
    leave-one-out sibling product is computed by prefix/suffix scans rather than
    by dividing out a message, since message entries are legitimately
    $\sim 10^{-16}$ in the tails.
\end{itemize}

\paragraph{EM on a tree.} \code{tree\_e\_step} produces the \textbf{same
$\Xis$ matrix} as the chain E-step: $\Xis[k,j]$ is the posterior mass on the
transition $u_j \to u_k$ summed over all tree edges. That is where the topology
stops being visible, so \emph{every} kernel in \code{src/kernels.py} consumes it
unmodified --- a chain-trained M-step and a tree-trained M-step are literally the
same code. The one genuinely new piece of bookkeeping is the evidence: messages
are renormalised at every node for numerical range, so $\log p(x)$ is recovered
by accumulating the discarded log-scales. That is what keeps the
monotone-ascent check available on trees.

\paragraph{Filtering.} \code{filter\_level=k} implements
Section~\ref{sec:filtering-def} faithfully; \code{filtered\_tree\_bp\_gaussian}
is the corresponding $\mathrm{BP}_k$.

\paragraph{A bimodal root.} \code{root\_separation}~$=\mu$ makes the root a
symmetric two-component mixture $\tfrac12\mathcal N(-\mu,1-\mu^2) +
\tfrac12\mathcal N(+\mu,1-\mu^2)$. Because this keeps unit variance, the leaf
covariance --- and therefore \emph{every eigenvalue and every predicted time} ---
is bit-identical to the Gaussian-root tree. Only the modality changes.
\code{tree\_root\_belief} returns the exact posterior of the root, which for a
two-component prior \emph{is} the class posterior.

\subsection{\code{src/spectral.py}}

The two time scales, in closed form and in this package's conventions:
the speciation crossover \code{speciation\_time}, the measurable
\code{commitment}, the AR(1) spectrum and its limit, the per-site excess entropy
and the collapse dataset size.

\subsection{Experiments}

\begin{center}\small
\begin{tabular}{@{}lll@{}}
\toprule
File & Parts & Subject \\
\midrule
\code{exp\_13\_speciation\_cascade.py} &
  \code{spectra}, \code{cascade}, \code{levels}, &
  the ladder, and what learns it \\
 & \code{ordering}, \code{block\_independent}, & \\
 & \code{speciation} & \\
\code{exp\_14\_memorization\_collapse.py} &
  \code{budget}, \code{collapse}, \code{time} &
  memorization and its time scale \\
\code{exp\_15\_bp\_oracles.py} &
  \code{ladder}, \code{alignment}, \code{mismatch} &
  the mismatched-oracle probe \\
\bottomrule
\end{tabular}
\end{center}

\subsection{Validation}

Every load-bearing object is computed twice, by independent routes:

\begin{center}
\begin{tabular}{@{}lll@{}}
\toprule
Quantity & Checked against & Agreement \\
\midrule
Analytic ultrametric spectrum & \code{eigh} of the dense covariance & $10^{-15}$ \\
Information-form tree BP & dense linear solve & $10^{-10}$ \\
Grid tree BP & information form & $2\times 10^{-6}$ \\
Filtered BP ($\mathrm{BP}_k$) & dense linear solve & $10^{-15}$ \\
Tree E-step evidence & closed-form Gaussian likelihood & rel.\ $10^{-6}$ \\
Root belief & dense linear solve & $10^{-5}$ \\
\code{commitment}$(t,\Lambda)$ & 40\,000 sampled trajectories & $<0.02$ \\
$\Xis$ mass & edge count & exact \\
EM monotonicity & --- & violation $=0$ \\
\bottomrule
\end{tabular}
\end{center}

\section{Derivations}
\label{sec:theory}

\subsection{The speciation crossover}
\label{sec:speciation}

Along an eigendirection of the clean covariance with eigenvalue $\Lambda$, the
variance-preserving forward process splits into signal and noise:
\begin{equation}
  \Var\bigl(\langle v, x_t\rangle\bigr) = \alpha_t^2\Lambda + \Delta_t,
  \qquad \alpha_t^2 = e^{-2t},\quad \Delta_t = 1-e^{-2t}.
\end{equation}
The mode is decided while signal dominates and undecided once noise does, so the
crossover sits at $\alpha_t^2\Lambda = \Delta_t$:
\begin{equation}
  \boxed{\;t_S(\Lambda) = \tfrac12\log(1+\Lambda)\;}
\end{equation}
which reduces to the familiar $\tfrac12\log\Lambda$ for large $\Lambda$. The
measurable form is the correlation between the projection at time $t$ and at
time $0$,
\begin{equation}
  \mathrm{commitment}(t,\Lambda)
  = \frac{\alpha_t\sqrt{\Lambda}}{\sqrt{\alpha_t^2\Lambda + \Delta_t}},
\end{equation}
which passes through $1/\sqrt2$ exactly at $t_S$. Verified analytically and
against 40\,000 sampled forward trajectories (agreement $<0.02$).

\begin{remark}[A distinction not to blur]
In P2, speciation is a \emph{symmetry breaking} and needs a multimodal
distribution to have something to choose between. A Gaussian prior has no class
to speciate into, so what \code{commitment} locates is an \emph{information
crossover} sharing the same criterion. Section~\ref{sec:symmetry} closes that
gap by making only the root bimodal, with the covariance held bit-identical.
\end{remark}

\subsection{The chain has a bounded spectrum; the tree does not}

For a stationary AR(1) chain with $\Cov = \rho^{|i-j|}$, the spectrum is bounded
at every length by the spectral density at zero frequency:
\begin{equation}
  \Lambda_{\max} \;\longrightarrow\; \frac{1+\rho}{1-\rho}.
\end{equation}
So $t_S$ \emph{saturates}. This is not a technicality: it says the chain family
studied in Layers 1--5 sits outside the regime P2 describes for image data,
where $\Lambda$ grows with dimension and $t_S$ diverges logarithmically.

For the balanced tree, writing $S_d$ for the sum of covariances over a depth-$d$
subtree,
\begin{equation}
  S_d = 1 + (b-1)\sum_{m=1}^{L-d} b^{m-1}\rho^{2m},
  \qquad
  \Lambda_d = S_{d+1} - b^{L-d-1}\rho^{2(L-d)},
\end{equation}
the latter with multiplicity $(b-1)b^{d}$, plus the uniform mode at $S_0$.
Exactly $L+1$ distinct eigenvalues, hence $L+1$ speciation times. When
$b\rho^2>1$ the top eigenvalue grows geometrically in depth, i.e.\
logarithmically in dimension.

\subsection{Filtering removes rungs from the ladder}
\label{sec:ladder-collapse}

This is where P1's knob and P2's time scale meet in one formula. Under the
correct filtering the leaf covariance is block-diagonal-plus-constant: the
depth-$(L-k)$ subtree covariance inside each of the $b^k$ blocks, and
$\rho^{2L}$ between them. Its eigenvectors are
\begin{enumerate}[nosep,label=(\roman*)]
  \item within-block contrasts --- the subtree's own non-uniform modes, each
    multiplicity multiplied by $b^k$;
  \item between-block contrasts, constant on each block and summing to zero
    across blocks, with the single eigenvalue
    $\Lambda_{\text{between}} = S_0^{(L-k)} - b^{L-k}\rho^{2L}$ and multiplicity
    $b^k-1$;
  \item the uniform mode,
    $S_0^{(L-k)} + (b^k-1)b^{L-k}\rho^{2L}$.
\end{enumerate}
So the top $k$ rungs \textbf{merge into one}, and the number of distinct
speciation times falls from $L+1$ to $L-k+2$, reaching 2 at $k=L$ where the
covariance is equicorrelated. Measured rung counts confirm this exactly
(Section~\ref{sec:exp15}).

\subsection{A closed-form collapse criterion}

For a Gaussian AR(1) chain the differential entropy is exactly extensive, so the
per-site excess entropy relative to the terminal measure $\mathcal N(0,1)$ is
\begin{equation}
  s = -\tfrac12\log(1-\rho^2),
\end{equation}
exactly, with no fitting, and the dataset size below which memorization is
forced is $e^{ns}$. At $\rho = 0.85$ that is $0.641$ nats per site. At finite
$t$ the noised law is $\mathcal N(0,\alpha_t^2 C + \Delta_t I)$, so
\begin{equation}
  s(t) = -\frac{1}{2n}\log\det\bigl(\alpha_t^2 C + \Delta_t I\bigr),
\end{equation}
and $n\,s(t_C) = \log N$ gives a \textbf{predicted collapse time with no fitted
constants}.

\section{Results}
\label{sec:results}

\subsection{The speciation ladder, measured}
\label{sec:cascade}

Depth-5 binary tree, $\rho = 0.9$, 32 leaves. Predicted $t_S$ against the
measured crossing of the commitment curve through $1/\sqrt2$, along the forward
process and along the reverse SDE driven by the exact tree-BP score:

\begin{center}
\begin{tabular}{@{}lrrrr@{}}
\toprule
Level & $\Lambda$ & $t_S$ predicted & forward & reverse SDE \\
\midrule
$-1$ (whole tree) & 14.271 & 1.363 & 1.368 & 1.336 \\
0 & 3.113 & 0.707 & 0.732 & 0.725 \\
1 & 1.804 & 0.516 & 0.521 & 0.520 \\
2 & 0.996 & 0.346 & 0.340 & 0.358 \\
3 & 0.498 & 0.202 & 0.208 & 0.222 \\
4 (sibling leaves) & 0.190 & 0.087 & 0.087 & 0.125 \\
\bottomrule
\end{tabular}
\end{center}

Six distinct transitions spanning $15.7\times$ in time \emph{within one
dataset}. Pearson $r = 0.9998$ for both columns; forward crossings agree to
$\leq 3.5\%$ at every level. The reverse column's deviations are
$-2, +3, +1, +4, +10, +43$ percent from coarsest to finest --- the monotone
growth toward the fine end identifies this as integrator error rather than a
failure of the prediction, since the finest level speciates at $t=0.087$, within
a factor of four of the sampler's $t_{\min} = 0.02$. Four of six levels agree to
$\leq 4\%$.

\paragraph{The chain has no ladder.} Same analysis on the AR(1) chain:

\begin{center}
\begin{tabular}{@{}lrrr@{}}
\toprule
$\rho$ & top eigenvalue, $n = 8 \to 512$ & limit & $t_S$ saturates at \\
\midrule
0.50 & $2.57 \to 3.00$ & 3.00 & 0.693 \\
0.85 & $5.49 \to 12.32$ & 12.33 & 1.295 \\
0.95 & $7.03 \to 38.52$ & 39.00 & 1.844 \\
\bottomrule
\end{tabular}
\end{center}

Sixty-four times the length moves the top mode by under $2\times$. A tree at
$\rho=0.9$ instead runs $3.12 \to 23.31$ over depths 2--6, its levels spanning
$t_S \in [0.087, 1.595]$. \textbf{This bounds how far Layers 1--5 generalise
toward image-like data}, and is the least comfortable finding of the session.

\subsection{Genuine symmetry breaking}
\label{sec:symmetry}

Depth-4 tree, $\rho=0.9$, $\mu=0.9$, 256 paths, $t_S$ predicted $1.136$. The
order parameter is the exact class posterior $P(\text{root}>0 \mid x_t)$, which
BP returns from the upward pass alone; using the sign of the tree-mean instead
would conflate the class with the within-class fluctuation.

\begin{center}
\begin{tabular}{@{}lrrr@{}}
\toprule
Root & correlation crossing & class-choice time & final magnetization \\
\midrule
two-component & 1.155 ($+1.7\%$) & \textbf{1.265} & 0.828 \\
Gaussian & 1.031 ($-9.2\%$) & 0.940 & 0.647 \\
\bottomrule
\end{tabular}
\end{center}

Two readings. First, the covariance-set crossing is where it should be for both
roots --- $1.155$ and $1.031$ bracket the predicted $1.136$, within the $\pm10\%$
the integrator costs at this resolution. \textbf{Adding modality does not move
it}, which is the control working. Second, the class choice for the
two-component root happens at $t=1.265$, \emph{earlier in the generation} than
the information crossover at $1.136$, and 35\% earlier than the corresponding
time for the Gaussian root: two well-separated classes can be told apart while
noise still swamps the within-class detail.

\begin{remark}[Precision, stated against my own result]
With 256 paths the binomial standard error on the agreement curve is $0.031$,
and the integrator contributes the $\pm10\%$ above. The $1.265$ versus $0.940$
gap is far outside both; the $+1.7\%$ agreement of the bimodal correlation
crossing is \emph{inside} them and must not be read as better than the $-9.2\%$.
\end{remark}

\subsection{Where each method's error lives in the hierarchy}
\label{sec:levels}

Depth-4 tree, $\rho = 0.9$, $N = 256$ trees, both network parameterizations
trained. EM recovered $\hat\rho = 0.9161$ (true $0.9$) and $\hat q = 0.1773$
(true $0.19$) in 150 iterations with monotonicity violation exactly 0.

\begin{center}
\begin{tabular}{@{}lrrr@{}}
\toprule
$t$ & EM-BP & best network & ratio \\
\midrule
0.1 & 0.0098 & 0.1852 ($\varepsilon$) & 18.9$\times$ \\
0.2 & 0.0125 & 0.2245 ($\varepsilon$) & 18.0$\times$ \\
0.4 & 0.0167 & 0.2489 ($x_0$) & 14.9$\times$ \\
0.8 & 0.0334 & 0.2587 ($x_0$) & 7.8$\times$ \\
1.6 & 0.0815 & 0.3332 ($x_0$) & 4.1$\times$ \\
\bottomrule
\end{tabular}
\end{center}

Grid tree BP given the true kernel scores $0.0000$ at every $t$, which is the
validation the rest of the part rests on.

More interesting is \emph{where} the error sits. The right null is a method
whose per-mode error is uniform, which would put a share of the total squared
error on each level proportional to its multiplicity ($1,1,2,4,8$ out of 16
here). Dividing the measured share by that null gives a concentration ratio,
whose spread across levels says how much a method's error knows about the
hierarchy:

\begin{center}
\begin{tabular}{@{}lr@{}}
\toprule
Method & spread of concentration ratio (min $\to$ max over $t$) \\
\midrule
network, $\varepsilon$-prediction & $\mathbf{1.6\times} - \mathbf{2.2\times}$ \\
network, $x_0$-prediction & $2.8\times - 12.6\times$ \\
EM-BP & $\mathbf{224\times} - \mathbf{6957\times}$ \\
\bottomrule
\end{tabular}
\end{center}

\textbf{The $\varepsilon$-network distributes its error uniformly over the
hierarchy at every noise level, as if the levels did not exist.} EM-BP's error
is concentrated by two to four orders of magnitude, and it \emph{migrates}: at
$t=0.1$ it sits on the finest levels, and by $t=1.6$ all of it is on the single
coarsest mode --- at large $t$ only the top mode survives the noise, so the
kernel misestimate has nowhere else to appear. The $x_0$-network is genuinely
intermediate, and is reported as such because it is the better parameterization
at large $t$.

\subsection{Memorization: the axis a BP score does not have}

The collapse mechanism is a statement about the \emph{empirical} score, whose
sufficient statistic is the training set itself: $N\times n$ numbers, growing in
both dataset size and dimension. A BP score's sufficient statistic is the fitted
kernel --- two numbers for a Gaussian AR(1) --- reached only through $\Xis$,
which is an average, and \emph{independent of $n$}. So the prediction is not
that BP memorizes less but that \textbf{the axis does not exist for it}.

Measured as nearest-training-chain distance of generated samples divided by the
same quantity for a fresh draw from the true prior ($1.0$ = no memorization,
$0$ = full collapse), $\rho = 0.85$:

\begin{center}
\begin{tabular}{@{}rrrrrrr@{}}
\toprule
$n$ & $N$ & wall $e^{ns}$ & empirical & DSM ($\varepsilon$) & EM-BP & true-prior BP \\
\midrule
8 & 64 & 88.8 & \textbf{0.457} & 0.800 & 1.079 & 1.076 \\
8 & 256 & 88.8 & \textbf{0.540} & 0.929 & 1.058 & 1.044 \\
8 & 1024 & 88.8 & \textbf{0.660} & 0.980 & 1.015 & 1.012 \\
16 & 64 & $1.5\times10^4$ & \textbf{0.332} & 0.751 & 0.972 & 0.992 \\
16 & 256 & $1.5\times10^4$ & \textbf{0.387} & 0.906 & 1.040 & 1.026 \\
16 & 1024 & $1.5\times10^4$ & \textbf{0.456} & 1.066 & 1.066 & 1.072 \\
32 & 64 & $4.3\times10^8$ & \textbf{0.268} & 0.790 & 1.021 & 1.024 \\
32 & 256 & $4.3\times10^8$ & \textbf{0.308} & 1.053 & 1.039 & 1.024 \\
32 & 1024 & $4.3\times10^8$ & \textbf{0.335} & 1.115 & 1.023 & 1.025 \\
\bottomrule
\end{tabular}
\end{center}

Both dependences come out in the predicted direction. At fixed $N$ the empirical
score memorizes \emph{more} as the chain gets longer ($0.457 \to 0.332 \to
0.268$ at $N=64$) --- the curse of dimensionality; at fixed $n$ it recovers as
$N$ grows. \textbf{EM-BP stays in $[0.97, 1.08]$ in every one of the nine
cells}, agreeing with the true-prior reference to $\leq 0.03$: its behaviour
does not depend on $n$ or $N$ at all.

\begin{remark}[The DSM column must be read with its standard deviation]
At $n=32$ the network's sample standard deviation is $1.19$--$1.27$, i.e.\
over-dispersed by 19--27\%, so its ratios above 1 there are inflated by
generating too broadly rather than by generalizing better. At $n=8$ and 16,
where its dispersion is closer to correct, its ratios sit below BP's. This is
why \code{sample\_std} and \code{lag1\_corr} are logged in every row.
\end{remark}

\paragraph{This is the mechanism behind Layer 5's headline.} The
$\geq 64\times$ data advantage and the $N^{-1/2}$ rate had been reported as
empirical facts. This is why they hold.

\subsection{The collapse time, predicted}

The criterion $n\,s(t_C) = \log N$, with no fitted constants, was tested over
\[
  n \in \{8,\,16,\,32\}
  \quad\times\quad
  N \in \{32,\,128,\,512,\,2048\}.
\]
\begin{itemize}[nosep]
  \item \textbf{Pearson $r = +0.990$} over the 9 settings where a finite-time
    collapse is predicted ($+0.995$ against the same measurement made on fresh
    rather than training chains).
  \item The relation is \textbf{affine, not proportional}:
    $\text{measured} = 1.716\times\text{predicted} + 0.276$, residual
    s.d.\ $0.029$. Forcing it through the origin gives residual s.d.\ $0.131$,
    $4.5\times$ worse.
  \item Independently, $\mathrm{d}t_C/\mathrm{d}\log N$ is $-0.039$ predicted
    versus $-0.061$ measured at $n=16$, and $-0.051$ versus $-0.079$ at $n=32$
    --- \emph{the same $1.55\times$ slope factor at both chain lengths}.
\end{itemize}

The criterion gets the ordering, both dependences, and the $N$-scaling to within
a constant factor; it does not get the absolute time. That is what should be
expected of a leading-order statement combined with an arbitrary ``half of
$\log N$'' landmark, and the offset is reported rather than absorbed into a
fitted constant.

\paragraph{Where it does not apply.} At $n=8$ with $N \geq 128$ the criterion
returns \emph{no finite-time collapse} (total excess entropy $4.487$ nats
against $\log 128 = 4.852$), yet collapse is measured at $t = 0.234, 0.170,
0.122$. There is no contradiction: the criterion is asymptotic in $n$, and a
finite system always collapses eventually as $t\to0$. Those three rows are
excluded from the correlation and listed rather than dropped.

\subsection{The mismatched-oracle probe}
\label{sec:exp15}

This is P1's actual instrument, and adopting it is the main methodological
result of the session.

\paragraph{The ladder collapses as derived.} Rung counts at $\rho=0.9$:

\begin{center}
\begin{tabular}{@{}lrrrrrrr@{}}
\toprule
& $k=0$ & 1 & 2 & 3 & 4 & 5 & 6 \\
\midrule
$L=4$: distinct rungs & 5 & 5 & 4 & 3 & 2 & --- & --- \\
$L=6$: distinct rungs & 7 & 7 & 6 & 5 & 4 & 3 & 2 \\
$L=6$: ladder span & 18.3 & 18.3 & 18.1 & 17.8 & 17.5 & 17.3 & 5.5 \\
\bottomrule
\end{tabular}
\end{center}

Exactly $L-k+2$, as derived in Section~\ref{sec:ladder-collapse}. That $k=0$ and
$k=1$ coincide independently reproduces P1's remark that those two cases share a
tree topology and differ only in the top transition probabilities.

\paragraph{The network implements the oracle matched to its training data.}
Train a denoiser on data filtered at $k_{\mathrm{train}}$, test on
\emph{unfiltered} data, and ask which $\mathrm{BP}_k$ its posterior mean is
closest to:

\begin{center}
\begin{tabular}{@{}lcr@{}}
\toprule
$k_{\mathrm{train}}$ & $\arg\min_k$ at $t = 0.1 / 0.2 / 0.4 / 0.8$ &
  margin over runner-up \\
\midrule
0 & 0 / 0 / 0 / 0 & --- \\
1 & 0 / 0 / 0 / 0 & \emph{exactly degenerate} \\
2 & \textbf{2 / 2 / 2 / 2} & $+1.8 / +4.1 / +4.3 / +1.6\ \%$ \\
3 & \textbf{3 / 3 / 3 / 3} & $+17.8 / +29.0 / +24.3 / +5.9\ \%$ \\
4 & \textbf{4 / 4 / 4 / 4} & $+63.3 / +60.4 / +35.3 / +6.2\ \%$ \\
\bottomrule
\end{tabular}
\end{center}

\textbf{16 of 16 cells} pick the oracle matched to the \emph{training}
distribution rather than the one matched to the test data, with the margin
growing in $k_{\mathrm{train}}$. This is P1's Figure 2 finding, transplanted to
diffusion. Two exclusions, both explained rather than hand-waved:

\begin{itemize}[nosep]
  \item \textbf{$k_{\mathrm{train}}=1$ is exactly degenerate with $0$ in a
    Gaussian tree} --- the distances agree to five decimals, gap $0.00\%$. This
    is a real limitation of the Gaussian analogue: P1's transition \emph{tensor}
    can correlate the two children given the parent, whereas a linear-Gaussian
    edge makes siblings conditionally independent by construction. The one case
    where their filtering alters probabilities without altering topology is
    \emph{unrepresentable here}.
  \item \textbf{At $t=1.6$ the $\arg\min$ stops measuring alignment} --- not
    because the oracles converge (their spread is still $0.22$) but because the
    network's distance to \emph{every} oracle is $\sim1.50$, seven times that
    spread. Every row carries \code{oracle\_spread} so this is checkable.
\end{itemize}

\subsection{The replicated architecture comparison}

A job pending from before the session completed during it, and it moves a
headline. Four seeds, oracle over both receptive-field radius and
parameterization for both networks, mean relative score error over five noise
levels, $\pm$ standard error over seeds:

\begin{center}
\begin{tabular}{@{}rrrr@{}}
\toprule
$N$ & EM-BP (13 params) & global MLP (25{,}505) & local CNN, oracle $r$ (6{,}081) \\
\midrule
128 & $\mathbf{0.0368 \pm 0.0067}$ & $0.3595 \pm 0.0044$ & $0.0774 \pm 0.0024$ \\
512 & $\mathbf{0.0174 \pm 0.0017}$ & $0.1794 \pm 0.0035$ & $0.0570 \pm 0.0007$ \\
2048 & $\mathbf{0.0123 \pm 0.0009}$ & $0.1246 \pm 0.0027$ & $0.0505 \pm 0.0015$ \\
\midrule
\multicolumn{4}{@{}l}{\emph{Ratio to EM-BP}} \\
128 & --- & $9.76 \pm 1.79$ & $2.10 \pm 0.39$ \\
512 & --- & $10.29 \pm 1.02$ & $3.27 \pm 0.32$ \\
2048 & --- & $10.15 \pm 0.75$ & $\mathbf{4.11 \pm 0.32}$ \\
\bottomrule
\end{tabular}
\end{center}

Two things, one in each direction. The structured architecture really does close
most of the gap: against the vanilla MLP the margin is $\sim10\times$ and flat
in $N$; against the locality-respecting CNN it is 2--4$\times$, so roughly 60\%
of the vanilla deficit was the architecture rather than the estimator.
\textbf{But the remaining margin grows with data rather than closing} ---
$2.10 \to 3.27 \to 4.11$, well outside the standard errors --- because EM-BP
keeps paying down parametric error (a factor 3.0 over $16\times$ data) while the
CNN flattens toward an approximation floor (a factor 1.5). ``More data will
close the gap'' is the natural guess and it is wrong here.

This supersedes the single-seed figures previously recorded ($3.3\times$ at
fixed $r=6$, $2.75\times$ with an oracle over $r$, both at $N=1024$).

\section{Corrections, and what came out negative}
\label{sec:corrections}

A research record is worth less than nothing if it only contains what worked.

\subsection{I had P1's filtering construction wrong}
\label{sec:correction-filtering}

Before reading the paper I implemented ``filter at level $k$'' as $b^k$
\textbf{independent} subtrees. P1 draws the depth-$k$ nodes conditionally
independently \emph{given the root}, so blocks remain correlated through it at
$\rho^{2L}$; independent blocks put a zero there instead. The implementation now
follows the paper, and a test pins the distinction rather than leaving it to a
comment. The old sweep survives, renamed \code{block\_independent} in code,
docs and output filenames, because it measures a harsher truncation that is
still well-defined --- but calling it filtering was a misattribution, not a
wording problem.

\subsection{I claimed an intersection was empty when it is not}
\label{sec:novelty}

An earlier version of the study note said the hierarchy-$\times$-diffusion
intersection was ``untouched by either paper''. Sclocchi, Favero and Wyart (PNAS
2025), \emph{cited by P1}, study exactly that and find a phase transition at
which reconstruction of high-level features drops sharply while low-level
features evolve smoothly. The claim is withdrawn. What the Gaussian setting adds
is more modest and more useful:
\begin{enumerate}[nosep]
  \item the whole ladder is analytic, so measurements are calibration against a
    known answer rather than discovery;
  \item filtering acts on the ladder in a computable way
    (Section~\ref{sec:ladder-collapse});
  \item their sharp/smooth split is \emph{reproduced} here --- smooth crossovers
    at every Gaussian level, genuine symmetry breaking only once the root is
    made bimodal, with the covariance held bit-identical
    (Section~\ref{sec:symmetry}).
\end{enumerate}

\subsection{The negative result: no sequential acquisition}

P1's headline dynamical finding does not reproduce here. Training on unfiltered
data and tracking which $\mathrm{BP}_k$ the denoiser is closest to, the
$\arg\min$ does move in the predicted direction at two of five noise levels ---
$2\to2\to0$ at $t=0.4$ and $2\to1\to1$ at $t=0.8$ over the first $\sim$1000
gradient steps --- and then \textbf{saturates for the remaining 31{,}000}. At
$t=0.1$, $0.2$ and $1.6$ it never leaves $k=0$. Nothing like their staircase.

Three plausible reasons, none verified: their oracles are far more
distinguishable (a discrete non-overlapping tensor versus a Gaussian edge --- see
the $k{=}0$/$k{=}1$ degeneracy); a denoising regression loss converges in far
fewer steps than MLM, which needs $\sim10^3$ epochs; and the architecture and
task differ. The honest summary is that \textbf{this setting is probably too easy
to resolve the phenomenon}, which is a different statement from the phenomenon
being absent.

The earlier probe (\code{exp\_13 ordering}: resolve the error onto levels and
watch it move) found nothing at all, and is retained only as a record of a
measurement that could not have answered the question.

\subsection{A non-monotonicity that is not a bug}

EM's \emph{error against the truth} is non-monotone in iteration count even
though its log-likelihood is monotone: total relative error $0.0729 \to 0.0083
\to 0.0162$ at 32/64/128 iterations, with $\hat\rho = 0.8708 \to 0.9068 \to
0.9158$ against a true $0.9$. Nothing is wrong. EM converges to the
\emph{maximum likelihood estimate}, which at $N=256$ is $\hat\rho = 0.9161$, not
to the true parameter; the trajectory merely passes close to $0.9$ on the way.
Reading the minimum at 64 iterations as ``the right stopping point'' would be
fitting the answer. The monotonicity violation is exactly 0 throughout, which is
the quantity that would actually signal a broken M-step.

\subsection{Tree EM converges, but slowly enough to look broken}

$\hat\rho = 0.7360 \to 0.7483 \to 0.7487$ at 50/100/150 iterations (true
$0.75$), with zero monotonicity violation throughout, and $\hat q = 0.4418$
against $0.4375$. Internal nodes are never observed, so the missing-information
fraction is far larger than on a chain and EM's linear rate is correspondingly
slower. An estimate read at 40 iterations is off by 0.06 and reads exactly like
a broken M-step. This cost real time to diagnose and is now asserted as a
\emph{rate} property, at two budgets, in the test suite.

\subsection{A confounded sweep, reported as confounded}

The \code{block\_independent} sweep shows the EM-BP advantage running
29--99$\times$ and not shrinking as correlations reach further, while the
network's error is flat in $k$ ($0.414 \to 0.347$). But fixing the number of
\emph{sequences} does not fix the number of \emph{edges} (16, 24, 28, 30 for
$k=1\dots4$), so larger $k$ hands EM more data at the same nominal budget. The
sweep moves two things at once, so the \emph{rising} advantage is not
attributable to correlation range. Only the negative is safe.

\section{Limitations and what should happen next}
\label{sec:limits}

\subsection{Limitations}

\begin{enumerate}[nosep]
  \item \textbf{P2 is unread.} Nothing here depends on its notation, but a
    write-up that cites it must use its symbols, and that correspondence has not
    been checked against the source.
  \item \textbf{The Gaussian tree cannot represent sibling correlation given the
    parent}, so P1's $k{=}0$ versus $k{=}1$ distinction is unrepresentable. A
    discrete-tensor tree would be needed.
  \item \textbf{The chain family is outside P2's regime.} Its top eigenvalue is
    bounded, so nothing about diverging speciation times transfers from Layers
    1--5 to image-like data without crossing that gap explicitly.
  \item \textbf{The collapse result uses a two-parameter kernel.} The obvious
    objection is that such a kernel \emph{cannot} memorize, so the result is
    trivial. \code{exp\_14 --set include\_mdn=True} runs the same comparison
    with a mixture-density kernel of several hundred parameters and the same
    $n$-independent statistic; it has not been run.
  \item \textbf{Inference cost remains the honest loss} of the whole programme:
    grid BP is $\mathcal O(nM^2)$ per evaluation against a few matmuls for a
    forward pass, and reverse diffusion calls the denoiser at every step.
\end{enumerate}

\subsection{Next, in priority order}

\begin{enumerate}[nosep]
  \item Read P2, and read Sclocchi/Favero/Wyart --- the closest prior work,
    whose territory an earlier version of my notes claimed was empty.
  \item The nonparametric collapse test (limitation 4), which separates the
    mechanism from the parameter count.
  \item A \textbf{discrete-alphabet tree}: closest to P1's actual data model,
    with exact vector messages and no closure. \code{src/discrete.py} does
    chains and \code{src/hierarchy.py} does continuous trees; the two have not
    been crossed. This would also make the $k{=}0$/$k{=}1$ case measurable.
  \item The speciation cascade under a \emph{learned} score rather than the
    exact one: which levels of the ladder each method gets right dynamically.
  \item A non-Gaussian tree. \code{tree\_bp\_grid} and \code{fit\_em\_tree}
    already support it with no new machinery.
\end{enumerate}

\section{Reproducibility}

\subsection{Running it}

\begin{verbatim}
cd research/nongaussian-bp
../../.venv/bin/python -m pytest tests/ -q      # 101 passed, ~3.3 min

python experiments/exp_13_speciation_cascade.py --list-parts
python experiments/exp_13_speciation_cascade.py --only cascade
python experiments/exp_14_memorization_collapse.py --only collapse
python experiments/exp_15_bp_oracles.py
\end{verbatim}

Add \code{--quick} to any experiment for a minutes-scale smoke run.
\code{--set KEY=VALUE} overrides any key printed in \code{params\_*.json};
unknown keys are rejected, so a typo in a job script fails immediately rather
than silently running defaults. \code{hpc/slurm\_layer6.sbatch} is an 8-task
SLURM array covering the parts worth a cluster.

\subsection{Where things are}

\begin{center}
\begin{tabular}{@{}ll@{}}
\toprule
Path & Contents \\
\midrule
\code{docs/PAPER\_CONNECTIONS.md} & the study note; \S0 records provenance \\
\code{docs/EM\_BP\_LEARNING\_COMPENDIUM.md} & all results, \S4.13 is Layer 6 \\
\code{docs/AGENT\_HANDOFF\_EM\_BP.md} & how to continue safely; 10 traps \\
\code{questions/QUESTIONS\_AND\_ANSWERS.md} & the question ledger, \S[P] is Layer 6 \\
\code{research/nongaussian-bp/src/hierarchy.py} & tree prior, tree BP, tree EM \\
\code{research/nongaussian-bp/src/spectral.py} & the two time scales \\
\code{research/nongaussian-bp/experiments/exp\_1\{3,4,5\}*.py} & the experiments \\
\code{research/nongaussian-bp/outputs/} & committed CSV + JSON + PNG \\
\code{research/nongaussian-bp/tests/test\_hierarchy.py} & 51 tests \\
\bottomrule
\end{tabular}
\end{center}

\subsection{Conventions that must not be broken}

Forward process $\mathrm dX = -X\,\mathrm dt + \sqrt2\,\mathrm dW$, so
$\alpha_t = e^{-t}$ and $\Delta_t = 1-e^{-2t}$. Score identity
$s(x,t) = -(x-\alpha_t\,\E[a\mid x])/\Delta_t$; every estimator goes through it,
which is why \code{identity\_residual} appears in every results CSV and must sit
at machine precision. All randomness flows through \code{src.utils.rng\_for},
which is reproducible \emph{across processes} --- it previously was not, and no
experiment in the package was bit-reproducible as a result. $\Xis$ uses
\code{K[out, in]} and sums to the edge count.

\section*{Summary of the session}
\addcontentsline{toc}{section}{Summary of the session}

\begin{enumerate}[nosep]
  \item Identified both papers as the advisors' own; read P1 in full through a
    server-side route after direct access was blocked.
  \item Built a hierarchical tree prior with a closed-form ultrametric spectrum,
    exact BP twice over, exact EM producing the same $\Xis$ as the chain, and
    the two diffusion time scales in closed form.
  \item Measured a ladder of six speciation transitions spanning $15.7\times$ in
    time, each within $3.5\%$ of prediction ($r = 0.9998$).
  \item Established that the AR(1) chain of Layers 1--5 has no such ladder,
    bounding how far the project generalises.
  \item Showed the memorization axis does not exist for a BP score, supplying
    the \emph{mechanism} behind Layer 5's headline sample-efficiency result.
  \item Verified a closed-form collapse time at $r = 0.990$, reporting the
    affine offset rather than absorbing it.
  \item Showed the $\varepsilon$-network's error is flat across hierarchy levels
    while EM-BP's is concentrated by 2--4 orders of magnitude.
  \item Reproduced P1's mismatched-oracle result in diffusion: 16 of 16 cells.
  \item Corrected two of my own claims --- the filtering construction and a false
    novelty claim --- and recorded one clean negative and one confound.
  \item Completed the replicated architecture comparison, which qualifies the
    Layer-5 headline and then partly un-qualifies it.
\end{enumerate}

Test suite: 50 $\to$ 101, all passing. \code{main} untouched at \code{cca0018}.

\end{document}
